\documentclass[11pt,a4paper]{article}
\usepackage[margin=1in]{geometry}
\usepackage{amsmath,amssymb,amsthm}
\usepackage{booktabs}
\usepackage{array}
\usepackage[T1]{fontenc}
\usepackage{lmodern}
\usepackage{microtype}
\usepackage{textcomp}
\usepackage{xcolor}
\definecolor{seclink}{HTML}{1F4E8C}
\usepackage[colorlinks=true,linkcolor=seclink,citecolor=seclink,urlcolor=blue]{hyperref}
\usepackage{parskip}
\usepackage[most]{tcolorbox}

\definecolor{ink}{HTML}{1A1A1A}
\definecolor{grey}{HTML}{5F5B54}
\definecolor{yellowbg}{HTML}{FBF3D5}
\definecolor{pinkbg}{HTML}{FADDE1}

\tcbset{plainbg/.style={boxrule=0pt, sharp corners, colframe=white,
  left=4mm, right=4mm, top=3mm, bottom=3mm, before skip=3mm, after skip=3mm}}
\newtcolorbox{sectionbox}{plainbg, colback=yellowbg}
\newtcolorbox{alarmbox}{plainbg, colback=pinkbg}

\newtheorem{claim}{Claim}
\newtheorem*{obs}{Observation}

\title{\vspace{-2em}The Trace and the Class}
\author{\normalsize Two obstructions, and the one thing each of them obstructs}
\date{\normalsize Claude, 20 September 2026}

\begin{document}
\maketitle
\vspace{-1em}

\begin{sectionbox}
\emph{A Sheaf Theory of Federation} closes by asking whether the trace obstruction
and the class in $H^1$ are the same obstruction, and offers two answers: they are
the same, or one of the two models is wrong.

This note argues for a third. They are different obstructions, neither model is
wrong, and they obstruct different properties of the same composite. Well-foundedness
governs \textbf{existence} --- is there an answer at all. The class governs
\textbf{canonicity} --- is the answer independent of the choice made to obtain it.
The statement relating them is not an equality but a composition: cutting an edge
trades divergence for choice, and $H^1$ is exactly the price of the choice.
\end{sectionbox}

\section{The question, as put}

From the closing box of \texttt{sheaf-federation.tex}:

\begin{quote}
The container note found that cycles obstruct a trace on $\triangleleft$. This note
finds that cycles carry a class in $H^1$. If those are the same obstruction, then
$H^1$ measures the failure of the trace to exist, and there ought to be a statement
relating them. If they are not the same, one of the two models is wrong about
federation and it would be useful to know which.
\end{quote}

The disjunction is too tight. Two obstructions can both be real, both be about
cycles, and still not be the same obstruction --- if they obstruct different things.
That is what happens here, and the program in this repository already contained the
evidence.

\section{Where each obstruction lives}

A container morphism is a pair $(u, f)$: $u$ covariant on prompts, $f$ contravariant
on responses. The two obstructions occupy different halves of that pair, and the
halves are independent by construction.

\emph{The Missing Trace} already localised one of them: prompts travel round the
cycle freely because $u$ is covariant and composes, and the failure is entirely in
$f$, where a response type ends up defined in terms of itself in negative position.
The sheaf note then found the class, and the class is carried in $f$ as well ---
\texttt{mobius.server.ts} applies $\varepsilon$ to the reply and nowhere else,
which is the only position descent allows it to occupy.

So both live in $f$, and it is natural to suspect they are one thing. The
observation that separates them is about $u$.

\begin{obs}
\texttt{mobius.server.ts} has four handlers that walk the ring. They differ only in
what stops them:
\emph{Ask} has fuel, which decreases and reaches zero;
\emph{Unroll} has fuel and a response type admitting \texttt{Exhausted};
\emph{Walk} has a visited list, and the node that sees itself returns the empty
product; \emph{Orient} has nothing --- no fuel, no visited list, no base case.
In all four, $\varepsilon$ appears in the same position, applied to the reply.
\end{obs}

The four differ in the \emph{forward} map's control flow, where no $\varepsilon$
occurs at all. Whatever distinguishes \emph{Orient} from \emph{Ask} is therefore
insensitive to the transition data. Writing the recursion out:
\[
  \mathrm{Orient}(i) \;=\; \varepsilon_i \cdot \mathrm{Orient}(i+1)
\]
there is no $i$ at which the right-hand side is not another call. The argument
never mentions $\varepsilon$; substituting $+1$ at every edge --- which is the
cylinder, the trivial class --- changes nothing about the recursion on the right.

\section{The 2$\times$2}

Running the ring confirms it. Each cell below was produced by the program in this
repository, on three live servers, with no code added to obtain it.

\begin{center}
\small
\begin{tabular}{llcccl}
\toprule
\textbf{class} & \textbf{repair} & $H^1$ & \textbf{terminates} & \textbf{canonical} & \textbf{witness} \\
\midrule
cylinder & none & $0$ & no  & ---  & abandoned after 8s, 757 hops \\
cylinder & cut  & $0$ & yes & yes  & all three cuts give $+,+,+$ \\
mobius   & none & $\neq 0$ & no  & ---  & abandoned after 6s, 660 hops \\
mobius   & cut  & $\neq 0$ & yes & no   & $+,-,-$ against $+,+,-$ against $+,+,+$ \\
\bottomrule
\end{tabular}
\end{center}

The load-bearing cell is the first. On the cylinder the twisted sheaf has two global
sections; there is nothing whatever obstructing descent; and the unbounded walk still
does not return. It was abandoned at eight seconds with the ring still turning. A
vanishing class did not buy termination, because termination was never what the class
was about.

The fourth cell is the other half. On the M\"obius band the cut repair
\emph{does} return --- every time, from every cut. What fails is that the three cuts
disagree, and any two of them differ by exactly the holonomy. So a non-trivial class
did not cost existence; it cost uniqueness.

\begin{sectionbox}
Well-foundedness is a property of $u$. The class is a property of $f$.
A container morphism is a pair, and the two components are chosen independently.
That is why the obstructions are independent, and it is not a coincidence about
this example.
\end{sectionbox}

\section{The statement relating them}

\begin{claim}
Let $g \in Z^1$ be transition data on the $n$-arc ring and let $\mathrm{Cut}_c(g)$ be
the section obtained by deleting edge $c$, rooting at $0$, and propagating. Then
\begin{enumerate}
\item $\mathrm{Cut}_c(g)$ is defined for every $c$ and every $g$. Existence does not
      see the class.
\item The unbounded walk is defined for no $g$, including $g = (+,+,+)$. Divergence
      does not see the class.
\item $\mathrm{Cut}_{c}(g) = \mathrm{Cut}_{c'}(g)$ for all $c, c'$ if and only if
      $[g] = 0$ in $H^1$.
\end{enumerate}
\end{claim}

Part (3) is the statement the closing box asked for, and it is not an equality
between the two obstructions. It says that once existence has been restored by
force, what remains is the class --- as the ambiguity of the repair, not as its
failure.

\texttt{independence.ts} checks all three on all eight elements of $(\mathbb{Z}/2)^3$
by enumeration rather than sampling, together with a fourth claim which is what makes
(3) a statement about $H^1$ at all:

\begin{claim}
Cut-agreement is gauge invariant. Each of the $24$ gauge transformations
$g \mapsto g \cdot d^0(a)$, $a$ supported at one node, moves the cocycle to a
different representative and never moves the verdict.
\end{claim}

Without this, (3) is a fact about chosen representatives. With it, it is a fact about
the class, and the two obstructions can be compared without a choice of coordinates
intervening.

\section{What follows for the repairs}

The guard chosen in $u$ decides which of the two failures one meets, and the four
handlers are four different bargains:

\begin{center}
\begin{tabular}{llll}
\toprule
\textbf{handler} & \textbf{guard} & \textbf{buys} & \textbf{costs} \\
\midrule
Orient & none & --- & does not return \\
Cut    & delete an edge & a section & the section depends on the cut \\
Ask    & fuel & a value & the value depends on the fuel's parity \\
Unroll & fuel, in the type & honesty & $F^n(\bot) = \bot$, for every $n$ \\
Walk   & visited & termination & nothing; it \emph{reports} the class \\
\bottomrule
\end{tabular}
\end{center}

\emph{Walk} is the interesting row. Guarding on the visited set restores
well-foundedness without forcing any choice, and the value it returns is the
holonomy itself. So a federation that cannot glue can still compute its own
obstruction, which is the practical form of the separation: the class is not an
error to be recovered from but a quantity to be reported.

\section{What this does not show}

\begin{enumerate}
\item \textbf{The budget is an instrument, not a proof.} A run abandoned at $757$
      hops has not demonstrated divergence. The proof of divergence is structural:
      the recursion has no base case, and the compiler says so independently in
      \texttt{mobius-type-tests.ts}, where \texttt{seqC(n2, ring)} is TS7022.
      Neither diagnostic mentions the signs, which is the actual argument.
\item \textbf{The coefficients are abelian and self-inverse.} $\mathbb{Z}/2$ cannot
      distinguish a loop from its inverse, so nothing here tests whether the
      separation survives non-abelian monodromy.
\item \textbf{The triple overlap is empty}, so the cocycle condition remains vacuous
      and this note inherits that limitation from the program it runs. The four-arc
      cover of \emph{Coherence} is where that is repaired, and the 2$\times$2 has
      not been rerun there.
\item \textbf{Only the JSV trace was ever ruled out.} Whether a weaker feedback
      notion --- Katis--Sabadini--Walters among them --- lands on containers remains
      unchecked, and this note does not check it. It does bear on it: a feedback
      category that hosted the trace would supply the operation, not the
      well-foundedness, so the first row of the 2$\times$2 would survive the
      upgrade unchanged.
\item \textbf{The ingredients were already here.} \texttt{mobius}'s own usage text
      says of the cut repair: ``ORIENT. Terminates. The answer depends on the cut.''
      What is new is not that sentence but the comparison it licenses --- that the
      same program diverges in the class where nothing is obstructed, and that the
      two facts put together answer a question that was left open.
\end{enumerate}

\begin{alarmbox}
\textbf{The revised question.} If well-foundedness and the class are independent,
a federation needs both to be discharged, and only one of them has an invariant.
$H^1$ says which twist, how much, around which loop. There is no comparable
invariant for the well-foundedness of a composite of container morphisms --- only
the binary fact that a guard is present or absent. Whether the well-founded ones
admit a grading as informative as $H^1$ is the question I would now put in place of
the one this note answers.
\end{alarmbox}

\end{document}
